\documentclass[../DoAn.tex]{subfiles}
\begin{document}

\begin{center}
    \Large{\textbf{TÓM TẮT NỘI DUNG ĐỒ ÁN}}\\
\end{center}
\vspace{1cm}

Các thiết bị di động hiện đại có khả năng thu nhận tín hiệu GNSS, dữ liệu cảm biến quán tính và luồng ảnh camera với tần suất cao. Tuy nhiên, trong quá trình theo dõi di chuyển trên đường, GNSS có thể yếu hoặc mất cập nhật do che khuất, đô thị dày đặc hoặc điều kiện thiết bị. Đồ án xây dựng một ứng dụng Android nhằm khảo sát cách dùng camera và cảm biến điện thoại để hỗ trợ cập nhật vị trí trực tiếp trên bản đồ khi GNSS suy giảm. Từ mục tiêu này, hệ thống được mở rộng thành ba bài toán: GNSS Viewer, Optical Flow View và LiveRouting kết hợp GNSS với Optical Flow/IMU. Ở GNSS Viewer, ứng dụng thu nhận vị trí, trạng thái vệ tinh và dữ liệu đo GNSS thô, sau đó ưu tiên phân giải vị trí vệ tinh theo Real GNSS PVT, IGS Broadcast Ephemeris, CelesTrak TLE/SGP4 và xấp xỉ từ azimuth/elevation. Kết quả được hiển thị trên bản đồ hai chiều, mô hình Trái Đất OpenGL 3D có ngày/đêm, Mặt Trăng, Mặt Trời, dữ liệu thời gian TAI-UTC của Orekit và màn hình AR phủ vệ tinh lên camera. Ở Optical Flow View, ứng dụng triển khai KLT và Farneback bằng OpenCV để phân tích camera thời gian thực, hỗ trợ ROI, object tracker, ghi video, lưu phiên đo và so sánh FPS, thời gian xử lý, số vector hoạt động cùng độ tin cậy. Với video upload, hệ thống hỗ trợ server FastAPI tự triển khai và backend ngoài của nhóm; pipeline RAFT dạng ONNX xử lý video theo job, upload/download chunk, hủy job, tải kết quả và dùng Cutie cho ROI đối tượng. Ở LiveRouting, ứng dụng theo dõi vị trí trên tuyến đường; khi GNSS không đủ tin cậy, camera assist được bật để Optical Flow và IMU hỗ trợ dead reckoning ngắn hạn, đồng thời vẽ các đoạn path phục vụ kiểm thử và quan sát lại.

\begin{flushright}
Sinh viên thực hiện\\
\begin{tabular}{@{}c@{}}
\textit{(Ký và ghi rõ họ tên)}
\end{tabular}
\end{flushright}

\end{document}
